\section{Introduction}
\label{sec:introduction}

Enterprise data platforms increasingly depend on pipelines that ingest, transform, validate, and deliver data across distributed systems. These pipelines support analytics, machine learning, regulatory reporting, customer-facing applications, and operational decision-making. Prior work has shown that data pipelines are integral to modern data-driven systems but remain vulnerable to quality, processing, and maintenance problems, including data-related defects, pipeline complexity, and hidden technical debt in production systems \cite{foidl2023pipelinequality,sculley2015hidden,xin2021production}. As organizations increase the number of sources, transformation stages, data contracts, quality checks, and downstream consumers, pipeline reliability becomes harder to maintain through manual monitoring and incident response alone.

Data-pipeline failures can arise from schema drift, missing-value spikes, duplicate records, freshness violations, unexpected volume changes, transformation defects, source-system failures, corrupted outputs, and unknown failure modes. These failures do not always appear as simple task crashes. A pipeline can complete successfully while producing stale, incomplete, duplicated, or semantically invalid data. As a result, pipeline reliability requires more than workflow execution status. It requires detection, diagnosis, recovery decisioning, validation, and evidence preservation.

Existing data-engineering systems address important parts of this problem. Workflow orchestration tools schedule and execute pipeline tasks, data-quality frameworks validate assumptions about data, observability tools expose anomalies and operational signals, incident-management workflows coordinate response, and policy-as-code systems provide guardrails for automated decisions \cite{apacheairflowdags2026,greatexpectations2026,dbt2026datatests,opa2026docs,googlesre2026monitoring}. Recent work on recurring data-pipeline validation and comparative evaluations of data-quality tools further reinforces the importance of automated monitoring and validation for production-style pipelines \cite{tu2023autovalidate,rehberger2026dqtools}.

However, these capabilities are often evaluated separately. Detection accuracy does not prove recovery safety. A successful retry does not prove that the underlying data issue was resolved. A remediation action does not prove that the recovered pipeline state is valid. Similarly, an alerting system may identify an incident without deciding whether the pipeline should retry, quarantine, fallback, stop, escalate, or continue. This creates a gap between observing pipeline failures and safely managing them.

This paper addresses that gap by presenting a policy-constrained self-healing control plane for data-pipeline failures. The framework evaluates self-healing as a staged lifecycle from failure detection to policy-approved and post-validated recovery. The design separates observability, diagnosis, remediation selection, policy approval, execution, validation, escalation, and evidence preservation so that the system is not rewarded for unsafe or unverifiable automation.

The central design principle is that automated recovery must be constrained by policy and validation. In production data environments, unrestricted remediation may create more risk than the original failure. For example, automatically dropping columns, rewriting outputs, backfilling stale data, suppressing failed quality checks, or bypassing validation may violate governance rules or corrupt downstream data products. Therefore, this framework counts recovery as successful only when the remediation is approved, executed, and verified. Failures that are detected and classified but cannot be safely repaired are escalated or contained rather than counted as successful self-healing.

The experimental evaluation uses a reproducible synthetic data-pipeline environment with controlled failure injection. The experiment includes 780 total trials: 690 injected failure trials, 60 healthy-control trials, and 30 boundary-control trials. Under the evaluated synthetic experiment configuration, the framework achieved 100\% failure detection accuracy and 100\% root-cause classification accuracy. Verified recovery was achieved in 52.17\% of injected failure trials. This separation between detection/classification performance and verified recovery performance is central to the paper's argument: self-healing data pipelines should be evaluated across multiple reliability stages, not only by whether a failure was detected or an automated action was attempted.

This paper makes the following contributions:

\begin{enumerate}
\item A staged self-healing control-plane architecture for data pipelines that separates failure detection, root-cause classification, remediation selection, policy approval, execution, validation, rollback, escalation, and evidence preservation.
\item A reproducible controlled failure-injection experiment with injected failure trials, healthy controls, and boundary controls.
\item A policy-constrained recovery model that prevents unsafe automation by requiring recovery actions to satisfy policy approval and post-recovery validation.
\item Separate reliability metrics for detection, diagnosis, recovery, runtime behavior, and classification quality, avoiding the misleading assumption that detection success implies recovery success.
\item A reproducibility package that preserves trial-level outputs, derived summaries, confusion-matrix data, figure-generation scripts, manuscript assets, and documentation for independent inspection and extension.
\end{enumerate}

The remainder of this paper is organized as follows. The Background and Related Work section positions the framework relative to workflow orchestration, data-quality validation, observability, AIOps, autonomic computing, production ML systems, and policy-governed automation. The Framework section describes the proposed self-healing control-plane architecture. The Methodology and Experimental Design sections describe the synthetic data generation, failure injection procedure, evaluation workflow, and reproducibility protocol. The Results section presents detection, classification, recovery, runtime, and confusion-matrix findings. The Discussion section explains the practical implications of the results. The Limitations and Threats to Validity section states the evaluation boundaries. The Conclusion summarizes the contribution and future research directions.
